\subsection{Validation with standard power-delay profiles}
\label{sec:standard_design}
An additional paired experiment uses the EPA and EVA power-delay profiles
specified in 3GPP TS~36.101, Annex~B.2.1, Tables~B.2.1-2 and B.2.1-3
\cite{etsi2010epa_eva}. All seven EPA and nine EVA paths are retained, with
the tabulated delays in nanoseconds and relative powers converted from dB
and normalized to unit total power. A sampling interval of 250~ns embeds
these profiles in the same 64-tone OFDM model: the subcarrier spacing is
62.5~kHz and the 16-sample CP lasts 4~$\mu$s. Maximum excess delays are
410~ns (1.64 samples) for EPA and 2510~ns (10.04 samples) for EVA. Both
satisfy the existing 12-sample memory bound and leave the four clean CP
pairs used by the calibration statistic. The standard power-delay profiles
are used with the paper's independent block-static Gaussian path gains;
the experiment does not implement the standardized Doppler processes or
an LTE-conformant waveform.

The protocol fixes EPA and EVA at SNRs 0, 10, and 20~dB before generating
three independent seed batches of 300 groups per cell (900 pooled,
5400 groups overall). Within each group, all methods use the same target
observations and the same 16-block noise-calibration window. Integer EB
and the continuous-delay receiver each use residual, CP, and known-noise
calibration; delay LS supplies a fixed baseline. Receiver parameters are
unchanged, including the six-component fitting cap, even though the physical
EPA/EVA profiles contain seven/nine paths. True profile parameters are used
to generate observations and evaluate error, but are not passed to the
practical estimators. The 2000 paired bootstrap resamples preserve each
seed batch's size. An independent finite-burst waveform check at
25-fold oversampling places every standard delay exactly on a 10~ns grid;
CP signal differences vanish numerically and the maximum FFT-response
error is $1.41\times10^{-15}$.
